\section{Discussion}\label{sec:discussion}

\paragraph{What the results mean for practice}
The central message is operational: a detector's usefulness as a \emph{filter} is a property of the pair (detector, calibration data), not of the detector alone. Our findings translate into a short deployment recipe.
(1) \emph{Never calibrate the miss budget on a convenient hold-out of the training distribution when the tool is going to be used on other projects or later releases.} The resulting number is not conservative or optimistic by a small margin; it was off by a factor of 2.7 on average, and by 3.7 in the worst case.
(2) \emph{Budget for labelled data from the deployment distribution.} About $100$ labelled vulnerable functions from the target (the new project, or the latest release window) are enough to recover the nominal miss rate on average; $25$--$50$ are usable with a conservative finite-sample correction; with fewer than $\lceil1/\alpha\rceil-1$ labelled positives no clearance is possible at all.
(3) \emph{Choose between the expected and the high-probability guarantee.} The expected guarantee is exact but a single calibration draw exceeds the budget by two points in 16--68\% of the draws at $k=100$; the PAC rule brings this to $\le7\%$ for roughly $30\%$ less automation.
(4) \emph{Calibrate per group when the group is observable.} Function length was an easy and effective stratifier; per-class calibration would require more labelled data than most teams have.
(5) \emph{Monitor.} Every cleared function that later turns out to be vulnerable is a labelled event; Algorithm~\ref{alg:triage} re-calibrates on such events.
(6) \emph{Check the economics.} If a missed vulnerability is worth hundreds of reviews, delegating a share of the code to a detector of this quality does not pay off; if it is worth tens, the saving of $27$--$32\%$ at $c_m=50$ on exchangeable data turns into a change of between $+5\%$ and $-27\%$ (a loss) under shift when the budget is calibrated on source data.

\paragraph{What the results mean for research}
We suggest that evaluations of vulnerability detectors report, besides AUROC/F1, the \emph{miss-versus-clearance} curve at fixed budgets on shift-aware splits, together with a leakage audit. AUROC correlated only weakly with the clearable share at $\alpha=5\%$ ($\rho=0.38$) and the optimism of in-distribution validation reached a factor of 3.8 for AUPRC. The audit also shows that cross-dataset claims between these two corpora must remove shared functions: otherwise the reported AUPRC is $3.5$--$5.6\times$ too high. Finally, the negative result for covariate-shift weighting is informative: weighting is the standard shift remedy in the conformal literature, and for source code its assumptions (a stable conditional $\Pr(Y|x)$ and a well-estimated density ratio in a very high-dimensional token space) do not hold well enough in practice.

\paragraph{Relation to trustworthy automated software engineering}
The special-issue themes of transparency, accountability and governance appear here as concrete, testable artefacts: the miss budget is an auditable statement ("at most $\alpha$ of the vulnerabilities in the calibration distribution are cleared"), its failure under shift is detectable once any labelled outcomes exist, the unevenness across sizes and CWEs shows where the statement hides risk, and the consistency study shows that the decision, not only the score, must be assessed under benign input changes. The approach transfers to other automated decisions in the development pipeline, e.g.\ filtering of software-composition-analysis alerts, triage of AI-generated code for review, or deciding which dependencies need manual vetting, wherever a rare, costly event is missed silently and labelled outcomes arrive over time.

\paragraph{Limitations that shape how to read the results}
The size dependence found by the explanation analysis suggests that part of the predictive signal in these corpora is a size confound rather than vulnerability semantics, which is a data issue and not a triage issue, but it limits what any detector trained on them can deliver. Stronger detectors would raise the achievable automation but would not change the logic of Propositions~\ref{prop:valid} and~\ref{prop:shift}. Online variants that adapt the threshold over time are a natural extension for drift that is gradual rather than abrupt.
